\section{模型评价与结论}
\subsection{方法的有效性与不足}
模型的主要作用是把可计算的相位信息与尚不确定的材料信息分开。色散坐标直接反映波数到传播相位的非线性关系，避免把相位折射率代入群光学周期；线性辅助参数与少量非线性参数分离，使背景与标定的处理可检查；两光束和多光束在相同数据条件下比较，减小了仅凭曲线形状判断机理的风险。分块预测、分角估计和折射率扰动又提供了不同层面的可靠性信息，防止把数值拟合精细误写成绝对测量准确。

主要不足来自未给定的材料与仪器条件。碳化硅使用$4H$普通光文献关系，无法识别本晶圆的实际取向、晶型和载流子贡献；硅使用固定外延层折射率及有效衬底响应，可能将真实的层内色散误差分配给其他参数。两角强度光谱未必在相同光斑取得，而共同厚度与共同响应的约束仍留下系统残差。有限波段内的条纹数量也限制了独立厚度信息量，不能将几千个采样点等同于几千次独立厚度测量。

后续测量宜补充晶型、温度、偏振及仪器分辨率，并在同一光斑重复测量；如果需要可追溯的绝对厚度，应增加独立折射率或另一种厚度测量作校准。模型方面可在这些信息具备后纳入各向异性传输、仪器卷积或厚度分布。继续增加无法被现有数据辨识的参数，未必能够提高结果可信度。

\subsection{各问题的回答}
\begin{table}[htbp]\centering
\caption{主结果及双角独立估计汇总}\label{tab:answer-summary}
\begin{tabular}{lrrrl}
\toprule
材料 & 共同厚度/$\mu$m & $10^\circ$独立/$\mu$m & $15^\circ$独立/$\mu$m & 主模型\\\midrule
SiC & 7.38431 & 7.40781 & 7.36078 & 色散两光束\\
Si & 3.40521 & 3.42055 & 3.39283 & 有效多光束\\
\bottomrule
\end{tabular}\end{table}

表\ref{tab:answer-summary}汇总当前光学假设下的共同厚度与分角估计，显示两种材料都需要同时保留角间一致性信息。
问题一的两束模型由同一出射波前上的光程差推得。相位差为$4\pi d_c\sigma\sqrt{n^2(\sigma)-\sin^2\theta}$再加界面反射相移；通过色散坐标之差可在未知绝对条纹级次时估计厚度。反射相移的常数部分改变明暗位置，随波数变化的部分则会影响厚度解释。

问题二在碳化硅$2000$--$3300\cmiv$窗口内计算，屏蔽$2300$--$2400\cmiv$局部异常，以文献普通光色散为条件，得到共同厚度$7.38431\um$，建议报告约$7.38\um$。两个角度分别为$7.40781\um$和$7.36078\um$，联合反射率RMSE为0.02440个百分点；这些数值可由随文代码和附件数据复现。

问题三通过往返振幅级数得到多光束模型。界面反射不为零、往返传播有足够振幅、出射光束在探测区域相干重叠，且仪器未平均掉相位信息，是产生可观测高阶结构需要检查的条件。硅光谱的多光束有效模型在相同参数数量下明显优于两束截断，共同厚度为$3.40521\um$，建议报告约$3.41\um$；其双角差异应同时保留。碳化硅主窗口的高阶形状校正仅改变厚度$-0.000158\um$，且未改善分块预测，所以不采用该修正。这个结论只适用于所检验的透明窗口，不排除其他波段存在多次反射。
